% Imporditud: tools/import_eksperimendid.py
% Allikas: Eesti füüsikaolümpiaadi eksperimendi ülesannete
%   kogu (Rael Kalda, 2023), ülesanne Ü51, lahendus L51.
\setAuthor{EFO žürii}
\setRound{piirkonnavoor}
\setYear{2008}
\setNumber{P E2}
\setDifficulty{2}
\setTopic{staatika}
\setVahendid{statiiv, statiiviklamber muhviga, joonlaud, 3 rahakummi, 2 kirjaklambrit, koormis massiga m = 55 g (mutter M20) ja tundmatu massiga keha (polt M14)}
\setPunktid{10}
\setAllikas{Eksperimendi kogu Ü51, lahendus lk 56}
\setLahendusseis{toores}

\prob{Polt}
Määrata tundmatu keha mass.

\hint

\solu
Kumminiidi pikenemine on võrdeline sellele rakendatud jõuga (2 p). Lõikame kum-

mirõnga lahti ja seome otstesse kirjaklambrid. Kinnitame kirjaklambri statiivi-

klambri külge. Mõõdame kirjaklambrite vahekauguse $h_0$. Riputame kirjaklamb-

ri otsa tuntud massiga keha ja mõõdame kirjaklambrite vahekauguse h. Leiame

kumminiidi pikkuse muutuse h $-$ $h_0$. Riputame kirjaklambri otsa tundmatu mas-

siga keha ja mõõdame kirjaklambrite vahekauguse H. Leiame kumminiidi pikku-

se muutuse H $-$ $h_0$. Kuna kumminiidi pikenemine on võrdeline sellele mõjuva

jõuga, seega on kumminiidi pikenemiste suhe võrdne selle otsa riputatud masside

suhtega. (3 p)

M = H $-$ $h_0$ $\Rightarrow$ M = H $-$ $h_0$ (1 p). m m h $-$ $h_0$ h $-$ $h_0$

Mõõtmised: Mõõtmised on dokumenteeritud ja nende põhjal on leitud H (2 p); tulemus on tõepärane (1 p). Mõõtmisi on korratud mitu korda (1 p).
\probend
